% paper/sections/08b_ccd_poisson.tex
%
% §8 圧力 Poisson 方程式の続き（08_pressure.tex からの分割）
% 内容：CCD-Poisson 演算子の行列構造・変密度拡張・Balanced-Force整合性・手法比較・単体検証
%
%% ═══════════════════════════════════════════════════════════════════════════════
\subsection{CCD-Poisson 演算子の行列構造と変密度拡張}
\label{sec:ccd_poisson_matrix}
%% ═══════════════════════════════════════════════════════════════════════════════

%% Modification I: O(h²) → O(h⁶) via CCD
%% 本節は §8.2 の目標を達成するための中核である：
%% (1) p, p', p'' を連立未知量として同時求解し O(h⁶) 精度を達成
%% (2) コロケート格子上のチェッカーボード不安定性を DCCD 散逸で抑制

\textbf{精度向上の動機：}
標準 $\Ord{h^2}$ 中心差分では界面近傍で「数値ノイズ」が生じ，
それが寄生流れ（parasitic currents）の主因となる．
CCD 法（第~\ref{sec:CCD}章）はこれを克服するため，
圧力 $p$，1次導関数 $p'$，2次導関数 $p''$ を\textbf{連立未知量として同時求解}する
（CCD ブロック系，付録~\ref{app:ccd_elliptic}）．
これにより速度補正に用いる圧力勾配の空間精度が $\Ord{h^6}$ に向上し，
長時間シミュレーションにおける離散化誤差の蓄積を根本的に低減する．

\textbf{コロケート格子のチェッカーボード（奇偶デカップリング）対策：}
コロケート配置では $p$ の奇数・偶数モードが PPE 左辺に現れず，
非物理的な格子スケール振動（checkerboard instability）が生じる．
本稿ではこれを Rhie--Chow 補間の代わりに
DCCD 散逸機構（§~\ref{sec:dccd_decoupling}，$\varepsilon_d=1/4$）で抑制する：
DCCD フィルタが PPE 右辺の $\bnabla_h\!\cdot\bu^*$ に対してチェッカーボードモードを除去し，
GFM ジャンプ補正との整合を保つ（式~\eqref{eq:gfm_ccd_div}）．

本節では，CCD 演算子行列 $\mathbf{L}_{\mathrm{CCD}}^{\rho}$（第~\ref{sec:CCD}章，
§~\ref{sec:ccd_elliptic}，式~\eqref{eq:L_CCD_2d_full}）を
PPE ソルバーとして明示的に活用するための実装上の詳細を述べる．

%% §8.2.1 (1D CCD-Laplacian 行列表現) および §8.2.2 (Neumann BC algbox) は
%% 付録 C.3（\ref{app:ccd_elliptic}，\ref{sec:ccd_neumann_bc}）に移動（CHK-086）．
%% ここでは要点のみ述べ，詳細は付録を参照する．

CCD-Laplacian 行列 $\mathbf{L}_{\mathrm{CCD}}^{(x)} = \mathbf{E}_2\,\mathcal{M}_{\mathrm{CCD}}^{-1}\,\mathcal{R} \in \mathbb{R}^{N \times N}$
の代数的導出（$\mathcal{R}$ の具体形，$\mathbf{E}_2$ 射影行列，スペクトル半径 $\approx 3.43/h^2$）は
付録~\ref{app:ccd_elliptic} に示す．
Neumann 境界条件（$p'_0 = 0$）の CCD ブロック Thomas 系への組み込み手順・変密度消滅・Dirichlet 比較表は
付録~\ref{sec:ccd_neumann_bc} に示す．

\subsubsection{2次元変密度 CCD-Poisson 演算子の完全定式化}
\label{sec:ccd_2d_full}

2次元変密度 PPE $\bnabla\cdot(1/\rho\,\bnabla p) = q$ の CCD 離散化を示す．
$x$ 方向の演算子（式~\eqref{eq:L_split}）を格子点 $(i,j)$ で評価すると：
\begin{align}
  \bigl(\mathcal{L}_{\mathrm{CCD}}^{\rho,x}\bm{p}\bigr)_{i,j}
  &= \frac{1}{\rho_{i,j}}\bigl(D_x^{(2)} p\bigr)_{i,j}
  - \frac{(D_x^{(1)}\rho)_{i,j}}{\rho_{i,j}^2}\bigl(D_x^{(1)} p\bigr)_{i,j}
  \label{eq:Lx_varrho_discrete}
\end{align}
$y$ 方向も同様に：
\begin{align}
  \bigl(\mathcal{L}_{\mathrm{CCD}}^{\rho,y}\bm{p}\bigr)_{i,j}
  &= \frac{1}{\rho_{i,j}}\bigl(D_y^{(2)} p\bigr)_{i,j}
  - \frac{(D_y^{(1)}\rho)_{i,j}}{\rho_{i,j}^2}\bigl(D_y^{(1)} p\bigr)_{i,j}
  \label{eq:Ly_varrho_discrete}
\end{align}

合計演算子：
\begin{equation}
  \bigl(\mathcal{L}_{\mathrm{CCD}}^{\rho}\bm{p}\bigr)_{i,j}
  = \bigl(\mathcal{L}_{\mathrm{CCD}}^{\rho,x}\bm{p}\bigr)_{i,j}
  + \bigl(\mathcal{L}_{\mathrm{CCD}}^{\rho,y}\bm{p}\bigr)_{i,j}
  \label{eq:L_CCD_2d_full}
\end{equation}

\noindent\textbf{大域行列形式（スパース直接法）：}
式~\eqref{eq:L_CCD_2d_full} の点ごと演算子は，クロネッカー積スパース行列
$\mathbf{L}_{\mathrm{CCD}}^{\rho}$（式~\eqref{eq:L_CCD_2d_kron}）として陽に組み立て，
スパース LU 法で直接求解できる（付録~\ref{app:ccd_kronecker}，\ref{app:ccd_lu_direct}）．
Matrix-Free 評価手順（各反復で Thomas 法を $x$・$y$ 方向に交互適用，$O(N_x N_y)$）は
付録~\ref{sec:ccd_matrix_free_eval} を参照のこと．

\subsubsection{GFM による界面処理と寄生流れの抑制}
\label{sec:ccd_balanced_force}

CSF 体積力モデルでは，表面張力項 $\sigma\kappa\bnabla\psi$ と PPE 左辺の離散化を
整合させる「Balanced-Force 条件」が必要であり，それでも CSF 近似誤差 $\Ord{\varepsilon^2} \approx \Ord{h^2}$
（$\varepsilon \sim h$）が寄生流れの物理的下限として残留した（CSF 時代の詳細は付録~\ref{app:csf_balanced_force} 参照）．

本稿では GFM（§~\ref{sec:gfm}）により界面を数学的不連続として扱う：
Young--Laplace ジャンプ $[p]_\Gamma = \kappa/We$ を PPE 右辺に直接組み込むことで
CSF 体積力を予測子ステップから完全に除去する．
これにより：
\begin{itemize}
  \item CSF 近似誤差 $\Ord{\varepsilon^2}$ による寄生流れの物理的下限が本質的に除去される．
  \item PPE 右辺は DCCD フィルタ済み発散 $+$ GFM ジャンプ補正（式~\eqref{eq:gfm_ccd_div}）のみから構成される
        （Rhie--Chow 補正は不要；§~\ref{sec:gfm}）．
  \item CCD 演算子が $p$・$p'$・$p''$ を同一スキームで評価するため，
        圧力勾配の離散化整合性が $\Ord{h^6}$ 精度で自動的に保たれる．
\end{itemize}

\begin{tcolorbox}[resultbox, title={GFM + CCD による寄生流れ抑制（FVM + CSF との比較）}]
\begin{center}
\renewcommand{\arraystretch}{1.2}
\begin{tabular}{lll}
\hline
手法 & 寄生流れの主因 & 理論的下限 \\
\hline
FVM + CSF + Rhie-Chow & 数値離散化誤差 $+$ CSF 物理誤差 & $\Ord{h^2}$（CSF 律速） \\
CCD + CSF + Balanced-Force & 数値成分低減，CSF 誤差残留 & $\Ord{\varepsilon^2} \approx \Ord{h^2}$（CSF 律速） \\
\textbf{CCD + GFM + DCCD}（本稿） & CSF 誤差を除去，数値成分 $\Ord{h^6}$ & \textbf{界面処理精度律速} \\
\hline
\end{tabular}
\end{center}
（表~\ref{tab:ccd_vs_fvm}，§~\ref{sec:balanced_force}，表~\ref{tab:accuracy_summary} も参照）
\end{tcolorbox}

\subsubsection{他手法との定量的比較}
\label{sec:ccd_vs_fvm}

\begin{table}[ht]
\centering
\caption{PPE 離散化手法の定量的比較（$N \times N$ 格子，$h = 1/N$）}
\label{tab:ccd_vs_fvm}
\renewcommand{\arraystretch}{1.4}
\begin{tabular}{lccccc}
\hline
手法 & 空間精度 & 記憶量 & 反復コスト & 寄生流れ ($N=64$) & 実装難易度 \\
\hline
FVM + BiCGSTAB & $\Ord{h^2}$ & $O(N^2)$ & $O(kN^2)$ & $O(10^{-2})$ & 低 \\
FVM + GMRES & $\Ord{h^2}$ & $O(N^2)$ & $O(kN^2)$ & $O(10^{-2})$ & 中 \\
CCD + ADI & $\Ord{h^6}$ & $O(N^2)$ & $O(N^2)$/iter & $O(10^{-5})$ & 中 \\
\textbf{CCD + 仮想時間陰解法} & $\bm{\Ord{h^6}}$ & $O(N^2)$ & $O(N^2)$/iter & $\bm{O(10^{-5})}$ & 中 \\
CCD + GMRES (matrix-free) & $\Ord{h^6}$ & $O(N^2)$ & $O(kN^2)$ & $O(10^{-5})$ & 高 \\
\hline
\end{tabular}

\smallskip
注：$k$ = 反復回数（BiCGSTAB/GMRES）；「寄生流れ」は静止液滴ベンチマーク
（§~\ref{sec:bench_droplet}）での $\|\bu_{\mathrm{para}}\|_\infty/(\sigma/\mu)$ の目安値．
\end{table}

%% 本章のまとめ（精度バランス表＋トレードオフ）は §8.4 の後に配置するため
%% 08h_pressure_summary.tex に移動（CHK-087: R5 — 読者が全節を読んでから結論に到達できるよう）．